% =====================================================================
%  scenarios.tex — Appendix C: Scenario Library
%  \input by facilitator_manual.tex
%
%  Ten ready-to-run scenarios. Each is engineered so that a DIFFERENT
%  control is the strongest answer. If every scenario rewarded the same
%  five cards, the game would teach one lesson repeatedly and students
%  would learn the answer rather than the reasoning.
%
%  All organizations are fictional. Never run a session against a real
%  client's systems.
% =====================================================================

\section{Scenario Library}
\label{app:scenarios}

Improvised scenarios are the most common reason a session goes flat, and
they also make cross-team comparison meaningless, since two teams arguing
about different situations are not disagreeing about anything. The ten
scenarios below are ready to read aloud.

Each carries a \textbf{teaching target}: the specific thing it is built to
surface. Each also carries an \textbf{expert key} naming which controls are
strong, partial, or weak \emph{for that situation}. The key is a scoring aid,
not an answer sheet --- a team that buys a ``partial'' card and defends it well
has done better work than one that buys a ``strong'' card without reasoning.

\begin{center}
\small
\begin{tabularx}{\textwidth}{@{}c X l c c@{}}
\toprule
\textbf{\#} & \textbf{Organization and teaching target} & \textbf{Budget} & \textbf{Tier} & \textbf{Care} \\
\midrule
1  & Food bank --- severity is consequence, not sophistication & 6 & Found. & \\
2  & Rural health clinic --- recovery beats prevention once it has happened & 6 & Found. & \\
3  & Domestic violence shelter --- when impact is physical safety & 6 & Found. & $\bullet$ \\
4  & Community land trust --- the boring control is the right one & 5 & Found. & \\
5  & School district --- you cannot buy your way out of a vendor's breach & 10 & Ext. & \\
6  & Municipal water utility --- the correct answer costs one point & 10 & Ext. & \\
7  & Legal aid office --- prevention fails, detection is the only lever & 10 & Ext. & $\bullet$ \\
8  & Public library --- you cannot protect what you do not know exists & 6 & Found. & \\
9  & Arts nonprofit --- what a grant paid for and nobody maintains & 6 & Camp. & \\
10 & Animal rescue --- the obligation that survives the incident & 10 & Ext. & \\
\bottomrule
\end{tabularx}
\end{center}

\vspace{0.4em}
{\small $\bullet$ Content that may land on a participant's personal experience.
Offer the redraw described in Section~7.4 \emph{before} reading it, addressed to
the team rather than to any individual.}

% --------- formatting macro ---------
\newcommand{\scen}[3]{%
  \vspace{1.0em}\noindent
  \textbf{#1. #2}\hfill{\small\textsc{#3}}\\[0.15em]
  \rule{\linewidth}{0.4pt}\\[0.25em]}
\newcommand{\sfld}[1]{\textit{#1}\hspace{0.4em}}

% =====================================================================
\scen{1}{Payroll Redirection at Harborview Food Bank}{Budget 6 \quad Foundational}
\sfld{Teaching target.} Severity is a property of who gets harmed, not of how
clever the attack was. The strongest control here is free and procedural.\\
\sfld{Read aloud.} Harborview Food Bank has three paid staff and about forty
volunteers. The bookkeeper works two days a week from home. Last Tuesday she
received an email that appeared to come from the board treasurer, in his usual
phrasing, asking her to update the payroll provider's bank details before
Friday's run. A follow-up voicemail in his voice confirmed it. Friday's payroll
is eighteen thousand dollars, roughly a month of operating cash.\\
\sfld{Suggested draw.} Method: Manipulation or Coercion $\cdot$ Impact: Financial
Wellbeing $\cdot$ Resources: Artificial Intelligence $\cdot$ Motivation:
Personal Gain\\
\sfld{Expert key.} \textbf{Strong:} Social Engineering Awareness (2), Deepfake
and AI Fraud Awareness (3), Email and Browser Protections (12). \textbf{Partial:}
Password and Authentication (5), Incident Response Planning (13).
\textbf{Weak here:} Encryption (14), Network Segmentation (18), Penetration
Testing (19) --- none of them touch a human being persuaded to type a new account
number.\\
\sfld{Debrief.}
\begin{itemize}
\item The attack was technically trivial and you rated it Very Severe. What made
it severe, and would it have been severe at a different organization?
\item Nobody can buy a callback rule. Which of your purchases would you trade
for one sentence in a written procedure?
\item The voice on the phone was cloned. Does awareness training still work when
familiarity stops being evidence?
\end{itemize}

% =====================================================================
\scen{2}{Ransomware at Cedar Ridge Health Clinic}{Budget 6 \quad Foundational}
\sfld{Teaching target.} Once an incident has happened, prevention cards are worth
nothing and recovery cards are worth everything. Teams that spent all six points
on prevention have to explain their Monday morning.\\
\sfld{Read aloud.} Cedar Ridge is a two-provider clinic serving about nine hundred
patients in a county with no other primary care. On Monday morning the front
desk cannot open the scheduling system. Every file on the shared drive has a new
extension, and a text file on the desktop asks for payment in cryptocurrency.
The practice manager thinks there is a backup, but it is on a drive plugged into
the same server, and nobody has ever restored from it.\\
\sfld{Suggested draw.} Method: Technological Attack $\cdot$ Impact: Physical
Wellbeing $\cdot$ Resources: Money $\cdot$ Motivation: Personal Gain\\
\sfld{Expert key.} \textbf{Strong:} Backup and Recovery (6) --- specifically the
isolated copy, since a backup on the same server is now encrypted too;
Incident Response Planning (13); Endpoint Detection and Response (16).
\textbf{Partial:} Software Updates (1), Email and Browser Protections (12),
Cyber Insurance (21). \textbf{Weak here:} Penetration Testing (19) --- an annual
report does not help on the morning it happens.\\
\sfld{Debrief.}
\begin{itemize}
\item The backup existed and was useless. What is the difference between having
a backup and having a recovery?
\item Impact is Physical Wellbeing, not Financial. What does a clinic without
scheduling actually do to a patient?
\item If you bought only prevention, describe your Monday.
\end{itemize}

% =====================================================================
\scen{3}{A Laptop Leaves the Building --- Wilder House}{Budget 6 \quad Foundational}
\sfld{Teaching target.} When Human Impact is physical safety, ordinary controls
carry stakes students have not previously attached to them.\\
\sfld{Handle with care.} This scenario involves a shelter and the safety of people
escaping violence. Offer the redraw to the team before reading it, and do not
require anyone to narrate.\\
\sfld{Read aloud.} Wilder House is a twelve-bed shelter. A caseworker's laptop was
taken from her car in a grocery store parking lot on Saturday. It holds an
offline copy of the intake spreadsheet: names, dates of birth, the schools
children attend, and current addresses for eleven families. The laptop has a
login password. Nothing else.\\
\sfld{Suggested draw.} Method: Physical Attack $\cdot$ Impact: Physical Wellbeing
$\cdot$ Resources: Entitlement $\cdot$ Motivation: Curiosity\\
\sfld{Expert key.} \textbf{Strong:} Encryption (14) --- full-disk encryption makes
this a hardware loss rather than a disclosure; Physical Security (8); Identity
and Access Management (4), on why an offline copy existed at all.
\textbf{Partial:} Asset Inventory (9), Incident Response Planning (13).
\textbf{Weak here:} Audit Logging (15), Network Segmentation (18) --- the device
is gone and off the network.\\
\sfld{Debrief.}
\begin{itemize}
\item A login password and full-disk encryption feel similar to a non-specialist.
What is the actual difference to whoever has the laptop now?
\item The spreadsheet was a convenience copy. Which control stops a convenience
copy from existing, and what does the caseworker lose when it does?
\item Who has to be told, how quickly, and what do you say to eleven families?
\end{itemize}

% =====================================================================
\scen{4}{The Volunteer Who Left in 2021 --- Fairmount Land Trust}{Budget 5 \quad Foundational}
\sfld{Teaching target.} The strongest answer is the least interesting card in the
deck. Teams reliably overlook it and buy something exciting instead.\\
\sfld{Read aloud.} Fairmount Land Trust has two staff and a rotating pool of
volunteers. A volunteer who managed the donation platform in 2021 moved out of
state and stopped responding. Last month a two-hundred-dollar test transaction
appeared, then reversed. The executive director does not know how many people
still have logins, and the platform bills annually to a credit card nobody can
locate the statement for.\\
\sfld{Suggested draw.} Method: Processes $\cdot$ Impact:
Financial Wellbeing $\cdot$ Resources: Entitlement $\cdot$ Motivation: Personal
Gain\\
\sfld{Expert key.} \textbf{Strong:} Account Lifecycle Management (7), Identity and
Access Management (4). \textbf{Partial:} Asset Inventory (9), Software Inventory
(10), Audit Logging (15). \textbf{Weak here:} Social Engineering Awareness (2)
--- nobody was tricked; a door was left open.\\
\sfld{Debrief.}
\begin{itemize}
\item Nobody attacked this organization. Is this a security incident?
\item The right answer costs one point and takes an afternoon. Why is it the one
nobody does?
\item How would this organization even produce a list of who has access?
\end{itemize}

% =====================================================================
\scen{5}{The Vendor's Breach --- Northbridge School District}{Budget 10 \quad Extended}
\sfld{Teaching target.} Sometimes nothing you can buy internally helps, and the
only real control was written a year earlier in a contract.\\
\sfld{Read aloud.} Northbridge is a rural district of four schools. Their bus
routing and student scheduling run on a hosted platform used by two hundred
districts. On Thursday a reporter calls the superintendent for comment on a
breach at that vendor. The district has had no notification. The vendor's
support line has a recorded message. The data includes student names, home
addresses, and bus stop times.\\
\sfld{Suggested draw.} Method: Indirect Attack $\cdot$ Impact: Relationships
$\cdot$ Resources: Technical Expertise $\cdot$ Motivation: Personal Gain\\
\sfld{Expert key.} \textbf{Strong:} Third-Party and Contractual Controls (17),
Incident Response Planning (13). \textbf{Partial:} Asset Inventory (9) extended
to data location, Cyber Insurance (21). \textbf{Weak here:} nearly everything
else --- Endpoint Detection, Segmentation, Secure Configuration and Software
Updates all protect systems the district does not own.\\
\sfld{Debrief.}
\begin{itemize}
\item Half your deck is useless here. What does that tell you about where a small
organization's real risk lives?
\item The superintendent learned from a reporter. What clause would have changed
that, and when would it have had to be written?
\item Bus stop times plus home addresses. Which Human Impact card is this really?
\end{itemize}

% =====================================================================
\scen{6}{Debug Mode in Production --- Millbrook Water Authority}{Budget 10 \quad Extended}
\sfld{Teaching target.} A one-point Foundational card beats every expensive option
on the table. Teams with ten points often overspend here.\\
\sfld{Read aloud.} Millbrook Water Authority serves eleven thousand residents.
Their public site posts water quality reports and lets residents pay bills. A
contractor built it four years ago and moved on. A resident emails to say that
adding a parameter to a URL shows an administrative page listing every account
and a button labelled ``post notice.'' The page has no login.\\
\sfld{Suggested draw.} Method: Processes $\cdot$ Impact:
Relationships $\cdot$ Resources: Technical Expertise $\cdot$ Motivation: Politics\\
\sfld{Expert key.} \textbf{Strong:} Secure Configuration (11) --- debug mode in
production is a configuration failure; Identity and Access Management (4).
\textbf{Partial:} Penetration Testing (19) would have found it, at three points;
Software Inventory (10); Third-Party and Contractual (17), on the contractor who
left. \textbf{Weak here:} Social Engineering Awareness (2), Backup and Recovery
(6).\\
\sfld{Debrief.}
\begin{itemize}
\item You had ten points. What is the cheapest purchase that actually closes this?
\item Penetration testing would have found it. Was it worth three points here, and
what would have to be true first?
\item A utility can post public notices. What is the worst version of this that
does not involve stealing anything?
\end{itemize}

% =====================================================================
\scen{7}{The Trusted Paralegal --- Eastgate Legal Aid}{Budget 10 \quad Extended}
\sfld{Teaching target.} Every preventive control in the deck fails, because the
person had legitimate access. Detection is the only remaining lever.\\
\sfld{Handle with care.} Involves surveillance of a person by someone they trusted.
Offer the redraw.\\
\sfld{Read aloud.} Eastgate Legal Aid handles housing and immigration matters. A
paralegal of six years has legitimate access to the case system. A client
reports that her landlord seems to know the contents of a filing that has not
been served. The paralegal and the landlord attend the same church. There is no
evidence of an intrusion, and nothing in the system is broken.\\
\sfld{Suggested draw.} Method: Processes $\cdot$ Impact:
Emotional Wellbeing $\cdot$ Resources: Entitlement $\cdot$ Motivation: Desire or
Obsession\\
\sfld{Expert key.} \textbf{Strong:} Audit Logging and SIEM (15) --- who read which
file and when is the only question that can be answered; Identity and Access
Management (4) with Least Privilege (24); Managed Detection and Response (20) if
affordable. \textbf{Partial:} Incident Response Planning (13), Third-Party and
Contractual (17) on employment terms. \textbf{Weak here:} Email and Browser
Protections (12), Encryption (14), Software Updates (1) --- the access was
authorized.\\
\sfld{Debrief.}
\begin{itemize}
\item Which of your purchases would have stopped this? If none, what would you
tell the client you can offer instead?
\item Logging tells you afterward. Is a control that only works afterward worth
two of ten points?
\item Least Privilege has a cost that is not measured in points. What does the
paralegal lose, and is the trade right?
\end{itemize}

% =====================================================================
\scen{8}{The Camera Nobody Remembered --- Delridge Public Library}{Budget 6 \quad Foundational}
\sfld{Teaching target.} Asset inventory is a precondition for everything else.
Teams cannot defend what they have not enumerated.\\
\sfld{Read aloud.} Delridge Public Library has public wifi, twenty patron
computers, a self-checkout system, and four security cameras installed by a
since-dissolved vendor in 2016. The county IT contractor reports that traffic is
leaving the library network at three in the morning, and traces it to a device
nobody on staff can identify. The library's network is flat: everything is on
the same wifi, including the circulation desk.\\
\sfld{Suggested draw.} Method: Technological Attack $\cdot$ Impact: Relationships
$\cdot$ Resources: Technical Expertise $\cdot$ Motivation: Curiosity\\
\sfld{Expert key.} \textbf{Strong:} Asset Inventory (9), Network Segmentation (18)
if the budget reaches it, Secure Configuration (11) for the 2016 default
credentials. \textbf{Partial:} Software Inventory (10), Audit Logging (15).
\textbf{Weak here:} Social Engineering Awareness (2), Backup and Recovery (6).\\
\sfld{Debrief.}
\begin{itemize}
\item At six points, segmentation may be out of reach. What do you buy first, and
what do you tell the library board about the rest?
\item The camera has been there for a decade. Which control would have caught it
in year one, and how much does that control cost?
\item Patron borrowing records are legally protected in many states. Does that
change your ranking?
\end{itemize}

% =====================================================================
\scen{9}{What the Grant Paid For --- Ridgeway Arts Collective}{Budget 6 \quad Campaign}
\sfld{Teaching target.} Built for multi-engagement play. The right move is often to
defer, and the card that lets you defer is the one that makes the cost of
deferring visible.\\
\sfld{Read aloud.} Ridgeway Arts received a technology grant in 2019 that paid for
a donor database, a website, and a year of support. The grant ended. The database
runs a version the vendor stopped patching in 2022. Migrating costs about four
thousand dollars the collective does not have. The board meets quarterly. Nothing
has gone wrong yet.\\
\sfld{Suggested draw.} Method: Technological Attack $\cdot$ Impact: Financial
Wellbeing $\cdot$ Resources: Technical Expertise $\cdot$ Motivation: Personal
Gain\\
\sfld{Expert key.} \textbf{Strong:} Software Inventory (10), Deferred Investment
(28) played honestly with a stated exposure, Backup and Recovery (6).
\textbf{Partial:} Network Segmentation (18) to isolate the unsupported system,
Compensating Control (27), Risk Acceptance (26) with a named owner.
\textbf{Weak here:} Software Updates (1) --- there are no patches to apply, which
is the whole problem, and teams reliably buy it anyway.\\
\sfld{Debrief.}
\begin{itemize}
\item Someone bought Software Updates. What happens when the vendor has stopped
shipping them?
\item If you deferred, name the exposure and say for how long. Would you put that
in writing to the board?
\item Revisit this in a later session and drift the threat. Was deferring still
the right call?
\end{itemize}

% =====================================================================
\scen{10}{Donor Records in a Shared Mailbox --- Second Chance Animal Rescue}{Budget 10 \quad Extended}
\sfld{Teaching target.} Obligations survive the incident. Notification and
disclosure are not technical controls and cannot be bought after the fact.\\
\sfld{Read aloud.} Second Chance runs on donations and a shared mailbox that six
volunteers access with the same password. On Friday the mailbox began sending
donation appeals to the entire contact list from an address that is not theirs.
The mailbox holds four years of correspondence, including scanned cheques with
routing numbers and a spreadsheet of major donors.\\
\sfld{Suggested draw.} Method: Manipulation or Coercion $\cdot$ Impact: Financial
Wellbeing $\cdot$ Resources: Money $\cdot$ Motivation: Personal Gain\\
\sfld{Expert key.} \textbf{Strong:} Password and Authentication (5) --- shared
credentials plus no multi-factor is the whole story; Account Lifecycle
Management (7); Incident Response Planning (13) for the notification duty.
\textbf{Partial:} Email and Browser Protections (12), Audit Logging (15), Cyber
Insurance (21). \textbf{Weak here:} Network Segmentation (18), Endpoint Detection
(16) --- nothing was installed on any machine.\\
\sfld{Debrief.}
\begin{itemize}
\item Six people share one password because it is genuinely convenient. What do
you replace it with that they will actually use? (This is the Human Factors card
in disguise.)
\item Which of your purchases helps you tell four thousand donors what happened?
\item Insurance may cover the cost. Does it cover the donor who stops giving?
\end{itemize}

% =====================================================================
\vspace{1.2em}
\subsection*{Writing Your Own}

A scenario that works has four properties. It names an organization concrete
enough to picture. It states a consequence in human terms rather than technical
ones. It has \emph{more than one} defensible answer, or teams agree in ninety
seconds. And it has at least one card that looks obviously right and is not ---
that card is where the argument happens.

Test a draft by asking what a team would buy with it, then asking whether a
different team could reasonably buy something else. If the answer to the second
question is no, the scenario is a quiz question rather than a game.
